\section{Methodology: Image Segmentation and Preprocessing}

\subsection{Overview}
\begin{enumerate}
    \item Image normalization
    \item Image segmentation
    \item Patch labeling
    \item Original image reconstruction with labels
\end{enumerate}

The entire preprocessing pipeline depends on a predefined patch size, which is
determined prior to all processing steps. This parameter influences multiple
stages of the pipeline, including image resizing, segmentation, and dataset
size.

\subsection{Patch Size Configuration}
This subsection defines the global patch size used throughout the preprocessing
pipeline.

Before the preprocessing pipeline, a fixed patch size must be determined. The
patch size is a global parameter that influences multiple stages of the
pipeline, including image resizing, segmentation, and the total number of
samples.

The number of image-segmented patches depends on the patch size. If the patch
size is small, the number of patches will be large, which increases the dataset
size. Therefore, the patch size directly affects the feasibility and performance
of the ML and DL models.

[Needed:Why patch size's choice is important]


\subsection{Image Normalization}
This step ensures that all images are transformed into a consistent format
compatible with the predefined patch size.

In this step, all images are resized to match the predefined patch size. This
guarantees that subsequent segmentation can be performed consistently across all
samples.

\subsection{Image Segmentation}
This step divides normalized images into smaller patches based on the predefined
patch size.

Images are segmented into fixed-size patches according to the predefined patch
size. The segmentation process directly depends on the patch configuration
defined earlier.

\subsection{Patch Labeling}
This step assigns labels to each segmented image patch for supervised learning.

Each image patch is assigned a label according to the selected labeling
strategy. These labels are used as ground truth for training the ML and DL
models.

\subsection{Image Reconstruction}
This step maps patch-level predictions back to the original image space.

After labeling or prediction, the patches are reassembled to reconstruct the
original image with corresponding labels. This allows visualization and
evaluation at the image level.

\subsection{Labeling Strategy}
This subsection defines how labels are assigned to image patches.

We decided to use machine auto-labeling instead of manual labeling, even though
we are training with small samples. The reason is that the original image is
large, and the number of patches increases as the patch size decreases. One of
the subsections will prove the correctness of this decision. In other words,
machine auto-labeling will not influence the experimental result.